% GENERATED FILE. Edit canonical lessons, then run npm run generate:books.
% canonical-source-hash: fnv1a64:c7c6186e430ae796
% canonical-lessons: ML-C405-vadikkuka, ML-C405-yukti, ML-C405-vastuta, ML-C405-sammatikkuka, ML-C405-nigamanam

\chapter{Arguing a Point}
\label{ch:ml-arguing-a-point}
\hlchaptermodality{405}

\begin{chapteropening}
\textbf{By the end of this chapter:} \emph{I can say to argue, logic, a fact, to admit and a conclusion.}

Complete the last lesson of chapter 405: I can say to argue, logic, a fact, to admit and a conclusion.
\end{chapteropening}

\section[\texorpdfstring{\ml{വാദിക്കുക} (vādikkuka)}{വാദിക്കുക (vādikkuka)}]{\ml{വാദിക്കുക} (vādikkuka) — to plead a case; to argue}

\label{lesson:ML-C405-vadikkuka}



\begin{quote}
Before the new one: say the Malayalam for a principal, then the Malayalam for tuition.
\end{quote}

\subsection*{You'll want to know: \ml{വാദിക്കുക}}
\textbf{\ml{വാദിക്കുക}} — \emph{vādikkuka} — \textquotedblleft{}to plead a case; to argue\textquotedblright{}.

\begin{quote}
\ml{അവൾ} \ml{പുതിയ} \ml{നിയമത്തിന്} \ml{എതിരെ} \ml{വാദിച്ചു}.
\end{quote}

\begin{quote}
\emph{avaḷ putiya niyamattinŭ etire vādiccu.}
\end{quote}

\begin{quote}
She argued against the new law.
\end{quote}

Past: \ml{വാദിച്ചു}, \emph{vādiccu}.

For a cause: -\ml{ക്ക്} \ml{വേണ്ടി} \ml{വാദിക്കുക}; against it: -\ml{ന്} \ml{എതിരെ} \ml{വാദിക്കുക}.

\begin{cousinweb}
\textbf{\ml{വാദിക്കുക}} (\emph{vādikkuka}) is built from Sanskrit \emph{vāda}, \textquotedblleft{}speech\textquotedblright{} (CDIAL 11507), from the root \emph{vad}, \textquotedblleft{}to speak\textquotedblright{}, with the Malayalam verb ending -\ml{ിക്കുക}. To argue is to keep speaking for one side.
\end{cousinweb}

\begin{grammarlens}[title={What you've built}]
One more word for giving reasons and opinions.
\end{grammarlens}

\subsection*{Your turn}
\begin{itemize}
\raggedright
  \item \emph{Say it:} \emph{vādikkuka}
  \item \emph{Say it:} \emph{vādikkuka}, once more
  \item \emph{Say it:} \emph{avaḷ putiya niyamattinŭ etire vādiccu.} — She argued against the new law.
  \item \emph{Recall:} say the Malayalam for a principal, then the Malayalam for tuition, then say \emph{vādikkuka} again
\end{itemize}

\begin{tcolorbox}[breakable,colback=teal!4,colframe=teal!35!black,title={Before you move on}]
What does \ml{വാദിക്കുക} mean? (\textquotedblleft{}To plead a case; to argue\textquotedblright{}.) Say it once more.
\end{tcolorbox}
\section[\texorpdfstring{\ml{യുക്തി} (yukti)}{യുക്തി (yukti)}]{\ml{യുക്തി} (yukti) — logic; reasoning}

\label{lesson:ML-C405-yukti}



\begin{quote}
Before the new one: say the Malayalam for tuition, then the Malayalam for to plead a case.
\end{quote}

\subsection*{You'll want to know: \ml{യുക്തി}}
\textbf{\ml{യുക്തി}} — \emph{yukti} — \textquotedblleft{}logic; reasoning\textquotedblright{}.

\begin{quote}
\ml{അവന്റെ} \ml{വാദത്തിൽ} \ml{ഒരു} \ml{യുക്തിയും} \ml{ഇല്ല}.
\end{quote}

\begin{quote}
\emph{avanṟe vādattil oru yuktiyuṁ illa.}
\end{quote}

\begin{quote}
There is no logic at all in his argument.
\end{quote}

A bookish word for good sense in an argument.

\begin{cousinweb}
\textbf{\ml{യുക്തി}} (\emph{yukti}) is borrowed from Sanskrit \emph{yukti}, first \textquotedblleft{}union, joining\textquotedblright{} and \textquotedblleft{}application\textquotedblright{}, and later \textquotedblleft{}reasoning, argument\textquotedblright{} (CDIAL 10480; Monier-Williams). It comes from the root \emph{yuj}, \textquotedblleft{}to join\textquotedblright{}, the root of \ml{യോജിക്കുക}: reasoning joins one thought to the next.
\end{cousinweb}

\begin{grammarlens}[title={What you've built}]
One more word for giving reasons and opinions.
\end{grammarlens}

\subsection*{Your turn}
\begin{itemize}
\raggedright
  \item \emph{Say it:} \emph{yukti}
  \item \emph{Say it:} \emph{yukti}, once more
  \item \emph{Say it:} \emph{avanṟe vādattil oru yuktiyuṁ illa.} — There is no logic at all in his argument.
  \item \emph{Recall:} say the Malayalam for tuition, then the Malayalam for to plead a case, then say \emph{yukti} again
\end{itemize}

\begin{tcolorbox}[breakable,colback=teal!4,colframe=teal!35!black,title={Before you move on}]
What does \ml{യുക്തി} mean? (\textquotedblleft{}Logic; reasoning\textquotedblright{}.) Say it once more.
\end{tcolorbox}
\section[\texorpdfstring{\ml{വസ്തുത} (vastuta)}{വസ്തുത (vastuta)}]{\ml{വസ്തുത} (vastuta) — a fact}

\label{lesson:ML-C405-vastuta}



\begin{quote}
Before the new one: say the Malayalam for to plead a case, then the Malayalam for logic.
\end{quote}

\subsection*{You'll want to know: \ml{വസ്തുത}}
\textbf{\ml{വസ്തുത}} — \emph{vastuta} — \textquotedblleft{}a fact\textquotedblright{}.

\begin{quote}
\ml{ഇത്} \ml{ഒരു} \ml{വസ്തുതയാണ്}, \ml{അഭിപ്രായമല്ല}.
\end{quote}

\begin{quote}
\emph{itŭ oru vastutayāṇŭ, abhiprāyamalla.}
\end{quote}

\begin{quote}
This is a fact, not an opinion.
\end{quote}

Plural: \ml{വസ്തുതകൾ}, \emph{vastutakaḷ}.

Plural: \ml{വസ്തുതകൾ}, the facts of a case.

\begin{cousinweb}
\textbf{\ml{വസ്തുത}} (\emph{vastuta}) is built on Sanskrit \emph{vastu}, \textquotedblleft{}a thing that exists; an affair\textquotedblright{} (CDIAL 11449), with the abstract ending -\ml{ത}; Sanskrit itself has \emph{vastutā} (Monier-Williams). A fact is what a thing really is.
\end{cousinweb}

\begin{grammarlens}[title={What you've built}]
One more word for giving reasons and opinions.
\end{grammarlens}

\subsection*{Your turn}
\begin{itemize}
\raggedright
  \item \emph{Say it:} \emph{vastuta}
  \item \emph{Say it:} \emph{vastuta}, once more
  \item \emph{Say it:} \emph{itŭ oru vastutayāṇŭ, abhiprāyamalla.} — This is a fact, not an opinion.
  \item \emph{Recall:} say the Malayalam for to plead a case, then the Malayalam for logic, then say \emph{vastuta} again
\end{itemize}

\begin{tcolorbox}[breakable,colback=teal!4,colframe=teal!35!black,title={Before you move on}]
What does \ml{വസ്തുത} mean? (\textquotedblleft{}A fact\textquotedblright{}.) Say it once more.
\end{tcolorbox}
\section[\texorpdfstring{\ml{സമ്മതിക്കുക} (sammatikkuka)}{സമ്മതിക്കുക (sammatikkuka)}]{\ml{സമ്മതിക്കുക} (sammatikkuka) — to admit; to agree}

\label{lesson:ML-C405-sammatikkuka}



\begin{quote}
Before the new one: say the Malayalam for logic, then the Malayalam for a fact.
\end{quote}

\subsection*{You'll want to know: \ml{സമ്മതിക്കുക}}
\textbf{\ml{സമ്മതിക്കുക}} — \emph{sammatikkuka} — \textquotedblleft{}to admit; to agree\textquotedblright{}.

\begin{quote}
\ml{ഒടുവിൽ} \ml{അവൻ} \ml{തന്റെ} \ml{തെറ്റ്} \ml{സമ്മതിച്ചു}.
\end{quote}

\begin{quote}
\emph{oṭuvil avan tanṟe teṟṟŭ sammaticcu.}
\end{quote}

\begin{quote}
At last he admitted his mistake.
\end{quote}

Past: \ml{സമ്മതിച്ചു}, \emph{sammaticcu}.

It also means to give consent: \ml{അച്ഛൻ} \ml{സമ്മതിച്ചു}, Father agreed. The noun is \ml{സമ്മതം}, consent.

\begin{cousinweb}
\textbf{\ml{സമ്മതിക്കുക}} (\emph{sammatikkuka}) is built from Sanskrit \emph{saṁmata}, \textquotedblleft{}agreed, approved\textquotedblright{} (CDIAL 12973), with the Malayalam verb ending -\ml{ിക്കുക}. It is the verb of \ml{സമ്മതം}, consent.
\end{cousinweb}

\begin{grammarlens}[title={What you've built}]
One more word for giving reasons and opinions.
\end{grammarlens}

\subsection*{Your turn}
\begin{itemize}
\raggedright
  \item \emph{Say it:} \emph{sammatikkuka}
  \item \emph{Say it:} \emph{sammatikkuka}, once more
  \item \emph{Say it:} \emph{oṭuvil avan tanṟe teṟṟŭ sammaticcu.} — At last he admitted his mistake.
  \item \emph{Recall:} say the Malayalam for logic, then the Malayalam for a fact, then say \emph{sammatikkuka} again
\end{itemize}

\begin{tcolorbox}[breakable,colback=teal!4,colframe=teal!35!black,title={Before you move on}]
What does \ml{സമ്മതിക്കുക} mean? (\textquotedblleft{}To admit; to agree\textquotedblright{}.) Say it once more.
\end{tcolorbox}
\section[\texorpdfstring{\ml{നിഗമനം} (nigamanaṁ)}{നിഗമനം (nigamanaṁ)}]{\ml{നിഗമനം} (nigamanaṁ) — a conclusion}

\label{lesson:ML-C405-nigamanam}



\begin{quote}
Before the new one: say the Malayalam for a fact, then the Malayalam for to admit.
\end{quote}

\subsection*{You'll want to know: \ml{നിഗമനം}}
\textbf{\ml{നിഗമനം}} — \emph{nigamanaṁ} — \textquotedblleft{}a conclusion\textquotedblright{}.

\begin{quote}
\ml{എല്ലാം} \ml{കേട്ട} \ml{ശേഷം} \ml{അവർ} \ml{ഒരു} \ml{നിഗമനത്തിൽ} \ml{എത്തി}.
\end{quote}

\begin{quote}
\emph{ellāṁ kēṭṭa śēṣaṁ avar oru nigamanattil etti.}
\end{quote}

\begin{quote}
After hearing everything, they reached a conclusion.
\end{quote}

To reach a conclusion is \ml{നിഗമനത്തിൽ} \ml{എത്തുക}.

\begin{cousinweb}
\textbf{\ml{നിഗമനം}} (\emph{nigamanaṁ}) is borrowed from Sanskrit \emph{nigamana}, which in the old books of logic names the summing up of an argument, its conclusion (Monier-Williams).
\end{cousinweb}

\begin{grammarlens}[title={What you've built}]
One more word for giving reasons and opinions.
\end{grammarlens}

\subsection*{Your turn}
\begin{itemize}
\raggedright
  \item \emph{Say it:} \emph{nigamanaṁ}
  \item \emph{Say it:} \emph{nigamanaṁ}, once more
  \item \emph{Say it:} \emph{ellāṁ kēṭṭa śēṣaṁ avar oru nigamanattil etti.} — After hearing everything, they reached a conclusion.
  \item \emph{Recall:} say the Malayalam for a fact, then the Malayalam for to admit, then say \emph{nigamanaṁ} again
\end{itemize}

\begin{tcolorbox}[breakable,colback=teal!4,colframe=teal!35!black,title={Before you move on}]
What does \ml{നിഗമനം} mean? (\textquotedblleft{}A conclusion\textquotedblright{}.) Say it once more.
\end{tcolorbox}
